\documentclass[conference]{IEEEtran}
% ============================================================================
%  Informe IEEE (IEEEtran) en espanol.
%  Parcial 02 - Inteligencia Artificial - Universidad Distrital.
%  Carrera de DARwIn-OP con Aprendizaje por Refuerzo en Webots.
%
%  ESTADO: las cifras de resultados estan marcadas [PENDIENTE] hasta ejecutar
%  los experimentos en una maquina con Webots. NO se inventan resultados.
% ============================================================================
\usepackage[utf8]{inputenc}
\usepackage[spanish]{babel}
\usepackage{amsmath,amssymb}
\usepackage{graphicx}
\usepackage{booktabs}
\usepackage{url}
\usepackage{hyperref}

\begin{document}

\title{Aprendizaje por Refuerzo con \emph{Safe Learning} para una Carrera de
Robots Humanoides DARwIn-OP en Webots: Comparacion entre PPO, DQN,
Actor-Critic y un Agente Propio Actor-Critic con Descubrimiento Poblacional}

\author{\IEEEauthorblockN{[Autores del grupo]}
\IEEEauthorblockA{Universidad Distrital Francisco Jose de Caldas\\
Inteligencia Artificial, Semestre 2026-3\\
Parcial 02}}

\maketitle

\begin{abstract}
Se desarrolla un entorno de carreras para el robot humanoide DARwIn-OP en el
simulador Webots, sobre el que se entrenan y comparan cuatro agentes de
aprendizaje por refuerzo (RL): un agente propio basado en Actor-Critic con
\emph{descubrimiento poblacional} (una estrategia de evolucion aplicada
periodicamente al actor), y tres algoritmos de referencia implementados con la
libreria Stable-Baselines3: PPO, DQN y Actor-Critic por ventaja (A2C). El
entorno expone una interfaz compatible con Gymnasium y se parametriza en tamano
de pista, posiciones de salida y meta, y numero y posicion de obstaculos. Se
incorpora \emph{safe learning} mediante un \emph{shield} (filtro de acciones
anti-colision), limites de velocidad y de amplitud de marcha, y una funcion de
costo separada con enfoque de Proceso de Decision de Markov con Restricciones
(CMDP). Se disenan 23 escenarios y se evaluan con multiples semillas, midiendo
recompensa acumulada, tasa de exito, tiempo a meta y metricas de seguridad.
\textbf{[PENDIENTE: resumen cuantitativo de resultados tras ejecutar los
experimentos.]}
\end{abstract}

\begin{IEEEkeywords}
aprendizaje por refuerzo, DARwIn-OP, Webots, PPO, DQN, Actor-Critic, evolution
strategies, safe learning, CMDP.
\end{IEEEkeywords}

% ===========================================================================
\section{Introduccion y descripcion del problema}
% ===========================================================================
El robot humanoide DARwIn-OP es una plataforma abierta ampliamente usada en
investigacion y docencia, simulable en Webots~\cite{webots}. El objetivo de este
trabajo es construir un entorno de \emph{carreras} en el que varios agentes
aprenden a desplazar al robot desde una linea de salida hasta una meta,
detectando y evitando obstaculos, y compitiendo entre si; gana el robot que
llega primero.

El problema combina tres dificultades: (i) control de un humanoide de 20 grados
de libertad, (ii) navegacion con obstaculos y (iii) aprendizaje con
restricciones de seguridad. Para separar el control de bajo nivel del
aprendizaje de alto nivel, el equilibrio y la marcha se delegan a un generador
de marcha (\emph{gait}); el agente de RL solo decide acciones de alto nivel
(avanzar, girar, detenerse). Esta decision se justifica en la
Seccion~\ref{sec:impl}.

Formalmente, modelamos la tarea como un Proceso de Decision de Markov (MDP)
$(\mathcal{S}, \mathcal{A}, P, r, \gamma)$ con observacion continua de dimension
12 (geometria relativa a la meta, distancia y angulo al obstaculo mas cercano,
e inercia del torso) y accion discreta de cinco opciones. La recompensa premia
el progreso hacia la meta y penaliza colisiones, caidas, salidas de pista y
tiempo (Seccion~\ref{sec:impl}).

% ===========================================================================
\section{Algoritmos}
% ===========================================================================
Esta seccion resume la logica de cada algoritmo y por que puede funcionar en el
problema de la carrera.

\subsection{PPO (Proximal Policy Optimization)}
PPO~\cite{ppo} es un metodo de gradiente de politica \emph{on-policy} que
estabiliza la actualizacion limitando (\emph{clipping}) cuanto puede cambiar la
politica en cada paso. Maximiza el objetivo recortado
\begin{equation}
L^{\text{CLIP}}(\theta) = \hat{\mathbb{E}}_t\!\left[
\min\!\big(\rho_t(\theta)\hat{A}_t,\;
\text{clip}(\rho_t(\theta), 1-\epsilon, 1+\epsilon)\,\hat{A}_t\big)\right],
\end{equation}
donde $\rho_t(\theta) = \dfrac{\pi_\theta(a_t\mid s_t)}{\pi_{\theta_{\text{old}}}(a_t\mid s_t)}$
es el cociente de probabilidades, $\hat{A}_t$ la ventaja estimada y $\epsilon$ el
parametro de recorte. Su estabilidad y bajo coste de ajuste lo hacen un buen
candidato para control con observaciones continuas como el de la carrera.

\subsection{DQN (Deep Q-Network)}
DQN~\cite{dqn} aprende una funcion de valor-accion $Q(s,a)$ con una red neuronal
y es aplicable a espacios de accion \emph{discretos}, como el nuestro. Minimiza
el error de diferencia temporal (TD) respecto a una red objetivo:
\begin{equation}
L(\theta) = \mathbb{E}_{(s,a,r,s')\sim \mathcal{D}}\!\left[
\big(r + \gamma \max_{a'} Q_{\theta^-}(s',a') - Q_\theta(s,a)\big)^2\right],
\end{equation}
con $\mathcal{D}$ el \emph{replay buffer} y $\theta^-$ los pesos de la red
objetivo. La exploracion se controla con una politica $\epsilon$-greedy. Al ser
discreto el conjunto de acciones (cinco), DQN encaja de forma natural.

\subsection{Actor-Critic por ventaja (A2C)}
A2C es la variante sincrona del Actor-Critic asincrono (A3C)~\cite{a3c}. Un
\emph{actor} $\pi_\theta$ propone acciones y un \emph{critic} $V_\phi$ estima el
valor del estado. El actor se actualiza con el gradiente de politica ponderado
por la ventaja:
\begin{equation}
\nabla_\theta J(\theta) =
\mathbb{E}\!\left[\nabla_\theta \log \pi_\theta(a_t\mid s_t)\,\hat{A}_t\right],
\qquad \hat{A}_t = \hat{R}_t - V_\phi(s_t),
\end{equation}
donde $\hat{R}_t$ es el retorno a $n$ pasos. Un termino de entropia favorece la
exploracion. Es un metodo simple y eficiente para acciones discretas.

\subsection{Agente propio: Actor-Critic con descubrimiento poblacional}
\label{sec:propio}
El agente propio (del Taller~2) parte de un Actor-Critic a $n$ pasos con bonus
de entropia (identico en espiritu a A2C) y lo \emph{aumenta} con una fase de
\emph{descubrimiento poblacional} inspirada en Evolution Strategies
(ES)~\cite{es}. La motivacion es que, en Actor-Critic, el actor solo cambia
segun lo que el critic opina del comportamiento que el actor \emph{ya} tiene;
si el critic esta equivocado, el agente puede no descubrir formas radicalmente
distintas (y mejores) de recorrer la pista.

Cada \texttt{disc\_every} pasos se construye una poblacion de \texttt{pop}
variantes del actor mediante perturbaciones gaussianas antiteticas
($\theta + \sigma\,\varepsilon_k$ y $\theta - \sigma\,\varepsilon_k$). Cada
variante \emph{juega carreras reales} y su aptitud (\emph{fitness}) es el
retorno real medio del episodio; el critic no interviene en esa decision. El
actor se mueve hacia las mejores variantes segun su rango:
\begin{equation}
\theta \leftarrow \theta + \texttt{es\_step}\cdot\sigma \sum_{k} w_k\,\varepsilon_k,
\qquad w_k \in [-0.5,\,0.5],
\end{equation}
donde $w_k$ es el rango de la variante $k$ centrado en cero. La variante de
\emph{control} baraja los pesos $w_k$ (direccion aleatoria de igual magnitud),
lo que permite comprobar si el beneficio proviene de \emph{seleccionar por score
real} o solo de \emph{perturbar} la red. Es \textbf{implementacion propia}
(Seccion~\ref{sec:impl}). Puede funcionar en la carrera porque la seleccion por
tiempo/recompensa real favorece explicitamente llegar antes a la meta.

% ===========================================================================
\section{Safe learning}
% ===========================================================================
El enfoque de \emph{safe learning} busca que el robot aprenda sin causarse dano
ni causarlo a otros robots u obstaculos. Se implementan restricciones explicitas
en tres niveles:

\begin{enumerate}
  \item \textbf{Terminacion y penalizacion por eventos inseguros}: caida (umbral
  de \emph{roll}/\emph{pitch} o caida del torso), colision y salida de pista
  terminan el episodio y reciben penalizacion.
  \item \textbf{Shield (filtro de acciones)}: antes de ejecutar la accion del
  agente, si esta llevaria al robot hacia el obstaculo mas cercano dentro de un
  umbral de distancia y casi de frente, se sustituye por un giro seguro hacia el
  lado libre. Se aplica por igual a los cuatro agentes.
  \item \textbf{Limites de actuacion}: la amplitud de avance y de giro de la
  marcha se recortan a valores seguros (limite de velocidad y de torque efectivo
  de las articulaciones).
\end{enumerate}

Ademas, se mantiene una \textbf{funcion de costo separada} con enfoque de
CMDP~\cite{cmdp}: el costo $c_t$ (caidas, colisiones, salidas de pista) se
acumula aparte de la recompensa para no distorsionar el objetivo de los agentes
base, y se reporta como metrica de seguridad independiente (numero de caidas,
colisiones y violaciones). Formalmente, el CMDP busca
$\max_\pi \mathbb{E}[\sum_t \gamma^t r_t]$ sujeto a
$\mathbb{E}[\sum_t \gamma^t c_t] \le d$, con $d$ un presupuesto de costo.

% ===========================================================================
\section{Implementacion y ejecucion}
\label{sec:impl}
% ===========================================================================

\subsection{Entorno y decisiones de diseno}
El entorno \texttt{DarwinRaceEnv} implementa la interfaz de
Gymnasium~\cite{gymnasium} (\texttt{reset}, \texttt{step},
\texttt{observation\_space}, \texttt{action\_space}) y delega la fisica a un
nodo Supervisor de Webots mediante un canal TCP local. La observacion es un
vector continuo de 12 dimensiones (posicion relativa a la meta, distancia y
angulo al obstaculo mas cercano, e inercia del torso). La accion es
\texttt{Discrete(5)}: detener, avanzar, avanzar-rapido, girar-izquierda,
girar-derecha.

\textbf{Justificacion del espacio de acciones y el gait.} Entrenar el equilibrio
de un humanoide de 20 DOF desde cero es costoso y fragil. Por ello cada accion
se traduce a parametros de un generador de marcha (amplitud de avance y de
giro), y el equilibrio lo resuelve la marcha, no el agente. Mantener el espacio
\emph{discreto} permite usar DQN sin adaptaciones y comparar los cuatro
algoritmos sobre el mismo \texttt{action\_space}.

\textbf{Justificacion de la arquitectura de competencia.} Cada agente se
\emph{entrena por separado} en pistas identicas (lo que mantiene el MDP
estacionario y la comparacion justa), y la \emph{carrera final} se ejecuta con
los cuatro robots DARwIn-OP simultaneamente en el mismo mundo (solo evaluacion,
sin aprendizaje), registrando el orden de llegada y los tiempos, como exige el
enunciado. Adicionalmente se compara por tiempos de las corridas separadas.

\textbf{Funcion de recompensa.}
\begin{equation}
r_t = w_{\text{prog}}\,(d_{t-1}-d_t) - p_{\text{time}}
      - p_{\text{col}}\mathbb{1}_{\text{col}}
      - p_{\text{fall}}\mathbb{1}_{\text{fall}}
      - p_{\text{out}}\mathbb{1}_{\text{out}}
      + b_{\text{goal}}\mathbb{1}_{\text{goal}},
\end{equation}
donde $d_t$ es la distancia a la meta.

\subsection{Librerias vs. implementacion propia}
PPO, DQN y A2C usan la libreria Stable-Baselines3~\cite{sb3} (no son
implementacion propia). El agente Actor-Critic con descubrimiento poblacional es
\textbf{implementacion propia} en PyTorch (Seccion~\ref{sec:propio}).

\subsection{Hiperparametros}
La Tabla~\ref{tab:hyper} resume los hiperparametros por defecto. En los
experimentos se varian algunos de ellos (p.\ ej.\ \emph{learning rate} y
$\gamma$).

\begin{table}[t]
\caption{Hiperparametros por defecto de cada algoritmo.}
\label{tab:hyper}
\centering
\small
\begin{tabular}{@{}llll@{}}
\toprule
\textbf{PPO} & & \textbf{DQN} & \\
\midrule
learning rate & $3\!\times\!10^{-4}$ & learning rate & $1\!\times\!10^{-3}$ \\
n\_steps & 1024 & buffer & 50\,000 \\
batch size & 64 & batch size & 64 \\
n\_epochs & 10 & train\_freq & 4 \\
$\gamma$ & 0.99 & target update & 1000 \\
gae\_$\lambda$ & 0.95 & $\epsilon$ inicial/final & 1.0 / 0.05 \\
clip range & 0.2 & exploration frac. & 0.2 \\
\midrule
\textbf{A2C} & & \textbf{Propio (AC+ES)} & \\
\midrule
learning rate & $7\!\times\!10^{-4}$ & learning rate & $7\!\times\!10^{-4}$ \\
n\_steps & 20 & n\_steps & 20 \\
$\gamma$ & 0.99 & $\gamma$ & 0.99 \\
ent\_coef & 0.01 & $\beta$ (entropia) & 0.01 \\
vf\_coef & 0.5 & pop / $\sigma$ & 6 / 0.01 \\
max\_grad\_norm & 0.5 & disc\_every / es\_step & 40k / 1.0 \\
\bottomrule
\end{tabular}
\end{table}

\subsection{Exploracion y participacion en la carrera}
Cada algoritmo (a)~\emph{explora} el entorno durante el entrenamiento (muestreo
estocastico de acciones o $\epsilon$-greedy) y (b)~\emph{participa} en la carrera
en evaluacion determinista (sin exploracion).

% ===========================================================================
\section{Experimentos y comparacion}
% ===========================================================================
Se definieron 23 escenarios (Tabla~\ref{tab:exps}) variando tamano de pista,
ancho, posicion de salida, posicion de meta y densidad de obstaculos. Cada
escenario se ejecuta con los cuatro algoritmos y al menos tres semillas. Los CSV
registran los episodios de entrenamiento (para las curvas de aprendizaje) y la
evaluacion (recompensa acumulada, exito, tiempo a meta, caidas, colisiones,
violaciones).

\begin{table}[t]
\caption{Escenarios de experimentacion (resumen).}
\label{tab:exps}
\centering
\scriptsize
\begin{tabular}{@{}lll@{}}
\toprule
ID & Descripcion & Variacion principal \\
\midrule
A1--A3 & Pista corta/media/larga sin obstaculos & tamano (largo) \\
B1--B2 & Pista estrecha / ancha & tamano (ancho) \\
C1--C3 & Salida desplazada / adelantada & posicion inicial \\
D1--D3 & Meta desplazada / cercana & posicion de meta \\
E1--E5 & 1 a 6 obstaculos & densidad de obstaculos \\
F1--F4 & Combinaciones (largo+estrecho+obstaculos) & mixta \\
G1--G3 & Giros fuertes / zigzag & curvas \\
\bottomrule
\end{tabular}
\end{table}

\subsection{Resultados}
\textbf{[PENDIENTE: ejecutar los experimentos en una maquina con Webots y
rellenar.]} Los marcadores siguientes se completan con \texttt{analyze.py}:

\begin{itemize}
  \item Tabla de metricas promedio ($\pm$ desviacion estandar) por algoritmo:
  recompensa acumulada, tasa de exito, tiempo a meta, caidas y colisiones
  (de \texttt{results/summary\_metrics.csv}).
  \item Figura de curvas de aprendizaje por escenario
  (\texttt{figures/learning\_curves\_*.png}).
  \item Figura de barras de tasa de exito (\texttt{figures/success\_rate.png}).
  \item Figura de boxplots de recompensa (\texttt{figures/reward\_boxplot.png}).
  \item Tabla/figura del orden de llegada en la carrera simultanea
  (\texttt{results/race\_summary.csv}).
\end{itemize}

\begin{table}[t]
\caption{Metricas promedio por algoritmo (evaluacion). [PENDIENTE]}
\label{tab:results}
\centering
\small
\begin{tabular}{@{}lccc@{}}
\toprule
Algoritmo & Recompensa (media$\pm$std) & Tasa de exito & Tiempo a meta \\
\midrule
Propio (AC+ES) & [PENDIENTE] & [PENDIENTE] & [PENDIENTE] \\
PPO & [PENDIENTE] & [PENDIENTE] & [PENDIENTE] \\
DQN & [PENDIENTE] & [PENDIENTE] & [PENDIENTE] \\
A2C & [PENDIENTE] & [PENDIENTE] & [PENDIENTE] \\
\bottomrule
\end{tabular}
\end{table}

% ===========================================================================
\section{Conclusiones y resultados}
% ===========================================================================
\textbf{[PENDIENTE: basar SOLO en los resultados reales obtenidos.]} Una vez
ejecutados los experimentos se debe: (i)~identificar que algoritmo obtuvo mejor
recompensa y tasa de exito y menor tiempo a meta; (ii)~indicar en que escenarios
fallo cada uno (p.\ ej.\ pistas estrechas o alta densidad de obstaculos) y el
motivo probable (exploracion insuficiente, inestabilidad de la marcha,
presupuesto de pasos corto); y (iii)~dar recomendaciones finales. No se incluyen
conclusiones cuantitativas hasta disponer de datos reales.

% ===========================================================================
\begin{thebibliography}{9}
% Cada referencia indica la seccion del informe a la que corresponde.
\bibitem{webots} O. Michel, ``Webots: Professional Mobile Robot Simulation,''
Cyberbotics. Simulador y PROTO del DARwIn-OP (Secc.\ I, IV).
\url{https://cyberbotics.com}

\bibitem{ppo} J. Schulman, F. Wolski, P. Dhariwal, A. Radford, O. Klimov,
``Proximal Policy Optimization Algorithms,'' arXiv:1707.06347, 2017 (Secc.\ II-A).
\url{https://arxiv.org/abs/1707.06347}

\bibitem{dqn} V. Mnih et al., ``Human-level control through deep reinforcement
learning,'' \emph{Nature}, vol.\ 518, pp.\ 529--533, 2015 (Secc.\ II-B).
\url{https://doi.org/10.1038/nature14236}

\bibitem{a3c} V. Mnih et al., ``Asynchronous Methods for Deep Reinforcement
Learning,'' arXiv:1602.01783, 2016 (Secc.\ II-C).
\url{https://arxiv.org/abs/1602.01783}

\bibitem{es} T. Salimans, J. Ho, X. Chen, S. Sidor, I. Sutskever, ``Evolution
Strategies as a Scalable Alternative to Reinforcement Learning,''
arXiv:1703.03864, 2017 (Secc.\ II-D).
\url{https://arxiv.org/abs/1703.03864}

\bibitem{cmdp} E. Altman, \emph{Constrained Markov Decision Processes}, Chapman
\& Hall/CRC, 1999 (Secc.\ III). Pagina del editor:
\url{https://www.routledge.com/Constrained-Markov-Decision-Processes/Altman/p/book/9780849303821}.
PDF del autor: \url{http://www-sop.inria.fr/members/Eitan.Altman/TEMP/h.pdf}

\bibitem{gymnasium} M. Towers et al., ``Gymnasium: A Standard Interface for
Reinforcement Learning Environments,'' arXiv:2407.17032, 2024 (Secc.\ IV).
\url{https://arxiv.org/abs/2407.17032}

\bibitem{sb3} A. Raffin et al., ``Stable-Baselines3: Reliable Reinforcement
Learning Implementations,'' \emph{JMLR}, vol.\ 22, no.\ 268, 2021 (Secc.\ IV).
\url{https://jmlr.org/papers/volume22/20-1364/20-1364.pdf}
\end{thebibliography}

\end{document}
